\chapter{المعلومات الكلاسيكية، والديناميات القابلة للعكس، والمصفوفات}

\lectureoverview{تطلب منا ميكانيكا الكم أن نستدل خارج النطاق الذي تدرَّب فيه حدسنا العادي. لذلك نبدأ خطوة أسبق، بالمعلومات الكلاسيكية: البتات، والتهيئات، وقوانين التحديث القابلة للعكس. ويقودنا النقاش نفسه طبيعيًا إلى المتجهات والمصفوفات، اللتين ستوفران لغة البتات الكمية.}

\section{حدس تشكّل لعالم غير ملائم}

حدسنا الفيزيائي قديم. فقبل زمن طويل من أن يكتب أي أحد معادلة، كان على الحيوان الذي يطارد فريسة أن يقدّر السرعة والاتجاه النسبيين. فإذا كان أسد يلحق بغزال، فإن الأسد يستجيب لمسافة فاصلة تتغير، وبالتالي، على نحو ضمني، للسرعة النسبية. وقد فُضّلت الأجهزة العصبية التي أحسنت التعامل مع مثل هذه الأحكام؛ أما التي أساءت التعامل معها فلم تخلّف نسلًا كثيرًا. وقد تدرّب كثير مما نسميه الحس السليم بشأن الحركة في هذا العالم العادي.

ذلك الإرث مفيد، لكنه ضيق. فلا شيء في الحياة اليومية يتحرك بسرعة تقارب سرعة الضوء، لذلك لم يزوّدنا التطور بأي غريزة تجاه الزمكان النسبي. ويتعلم الفيزيائي أولًا الهندسة المجردة، ثم لا يكتسب حدسًا جديدًا لها إلا لاحقًا. وميكانيكا الكم أصعب من ذلك أيضًا. فصورتنا المألوفة للجسيم مستمدة من الحجارة والسهام: في كل لحظة يكون للجسم موضع، ويشكّل تعاقب المواضع مسارًا، وتعرّف تغيرات الموضع سرعة. لكن الإلكترون لا ينسجم مع تلك الصورة. ولا نظير لعدم اليقين في الموضع وغياب مسار كلاسيكي في الخبرة التي شكّلت حواسنا.

لا مفارقة في هذا الإخفاق في الحدس. ولا تكون هناك مفارقة إلا إذا كانت الغرائز المتكيفة مع الحجارة والحيوانات والسرعات المعتدلة تعرف سلفًا شيئًا عن الإلكترونات. والاستجابة الصحيحة هي أن نبني حدسًا مختلفًا انطلاقًا من مجموعة صغيرة من القواعد الدقيقة. ولن يكون التركيز هنا على المسار المعتاد عبر ميكانيكا الموجة ومعادلة Schr\"odinger. بل سنركز على البنية المنطقية لميكانيكا الكم وعلى المعلومات الكمية. وقبل أن نسأل ما البت الكمي، يجب أن نفهم على وجه الدقة ماذا يفعل البت العادي.

\section{البتات وكمية المعلومات}

الفيزياء لا تنفصل عن المعلومات. وأحيانًا تُعبَّر تلك المعلومات بعدد حقيقي، لكن أبسط قدر ممكن من المعلومات هو جواب عن سؤال له بديلان فقط: نعم أو لا، وجه أو ظهر، أعلى أو أسفل.

\begin{definition}
إن \emph{البت الكلاسيكي} هو سؤال فيزيائي له إجابتان ممكنتان بالضبط.
\end{definition}

يمكن أن نرمز إلى هاتين الإجابتين بـ \(0\) و\(1\). ولا تحمل هذه الرموز أي معنى ذاتي؛ فأي رمزين متميزين يفيان بالغرض. ولأن ترميز ديراك سيصبح بعد قليل لا غنى عنه، فإننا ندخله هنا مع أنه زائد عن الحاجة:
\begin{equation}
  \ket{0},
  \qquad
  \ket{1}.
\end{equation}
يشكّل الخط العمودي والقوس الزاوي \emph{كيت}، وهو النصف الأيمن من برا--كيت. ولم يدخل أي شيء كمي لمجرد أننا غيّرنا الترميز.

بالنسبة إلى عدة بتات قابلة للتمييز، اكتب الإجابات بترتيب ثابت. فست عملات معدنية قابلة للتمييز، مثلًا، قد تكون لها التهيئة
\begin{equation}
  \ket{0\,1\,1\,0\,1\,0}.
\end{equation}
وتكون الفواصل اختيارية متى فُهم الترتيب. للبت الواحد حالتان ممكنتان، وللبتين أربع حالات، وكل بت يضاف يضاعف العدد. ومن ثم فإن منظومة مؤلفة من \(n\) بتات تمتلك عددًا من الحالات مقداره
\begin{equation}
  N_s=2^n
  \label{eq:l01-state-count}
\end{equation}
وبالعكس،
\begin{equation}
  n=\log_2 N_s.
  \label{eq:l01-bits-from-states}
\end{equation}

تشير المعادلة~\eqref{eq:l01-bits-from-states} إلى تعريف مفيد حتى عندما لا يكون عدد الحالات قوة للعدد اثنين:
\begin{equation}
  I=\log_2 N_s.
  \label{eq:l01-information}
\end{equation}
ومن ثم يحمل نرد ذو ستة أوجه \(\log_2 6\) بتًا من المعلومات. وليس من الضروري أن يكون الجواب عددًا صحيحًا. وستكون الأنظمة المؤلفة حرفيًا من مفاتيح ثنائية أمثلة مريحة، لكن اللوغاريتم هو المقياس الأعم.

\section{ترميز عالم يبدو متصلًا}

لا تقتصر البتات على العملات المعدنية. فكل كمية مقاسة معروفة بدقة منتهية يمكن كتابتها على هيئة سلسلة ثنائية منتهية. وافترض، مثلًا، أن درجة حرارة غرفة لا تُطلب إلا إلى عدد مختار من الأرقام المعنوية. عبّر عن القيمة المقاسة بالأساس الثنائي واحتفظ بالخانات اللازمة لتلك الدقة. وتتطلب دقة أكبر سلسلة أطول، من غير أي تغيير في المبدأ.

ويمكن معالجة مجال يتغير مكانيًا بطريقة مماثلة. اقسم الغرفة إلى شبكة مرتبة دقيقة من الخلايا، وسجّل قيمة المجال في كل خلية بدقة منتهية، ثم صِل السلاسل الثنائية الناتجة بعضها ببعض. إن تصغير الخلايا يحسّن الاستبانة المكانية؛ والاحتفاظ بمزيد من الخانات الثنائية يحسّن استبانة قيم المجال. ويمكن ترميز مجالات درجة الحرارة، والمجالات الكهربائية، والمجالات المغناطيسية كلها بهذه الطريقة إلى أي دقة منتهية مطلوبة.

وتعطي تهيئات الجسيمات صورة ثنائية أشد حرفية. فقسّم الفضاء إلى خلايا صغيرة بما يكفي بحيث تحتوي كل واحدة منها على جسيم واحد على الأكثر. وعيّن
\begin{equation}
  n_i=
  \begin{cases}
    1,&\text{cell \(i\) is occupied},\\
    0,&\text{cell \(i\) is empty}.
  \end{cases}
  \label{eq:l01-occupancy-bit}
\end{equation}
والتهيئة الكاملة هي السلسلة المرتبة
\begin{equation}
  \ket{n_1n_2n_3\cdots}.
\end{equation}
وهي جواب طويل جدًا عن متتالية من أسئلة نعم أو لا: هل يوجد جسيم في الخلية الأولى، وفي الثانية، وهكذا؟ وهذه هي النسخة الفيزيائية للعبة \emph{عشرون سؤالًا}، إلا أن عدد الأسئلة قد يكون هائلًا.

يبين هذا البناء لماذا تستطيع الحواسيب الرقمية محاكاة الفيزياء. فالحاسوب يخزن سلاسل منتهية، بينما تحوّل الاستبانة المنتهية الكميات المقاسة والحقول إلى سلاسل من هذا النوع على وجه التحديد. ويمكن تحسين التقريب على نحو منهجي بتنقيح الشبكة، أو خطوة الزمن، أو عدد الخانات المحتفظ بها.

\begin{classroomqa}
\lecturetimestamp{00:29:43}
\audiencequestion{هل يمكن اختيار دقة تمثيل درجة الحرارة اعتباطيًا؟}
\lecturerresponse{تُختار وفقًا للتقريب الذي نريده. فالسلسلة المنتهية تعطي دقة منتهية؛ والاحتفاظ بعدد أكبر من الخانات الثنائية يعطي تقريبًا أدق. ولا يوجد هنا ادعاء بأن سلسلة منتهية تحدد عددًا حقيقيًا كلاسيكيًا اعتباطيًا تحديدًا تامًا.}
\end{classroomqa}

\begin{classroomqa}
\lecturetimestamp{00:30:11}
\audiencequestion{ماذا يحدث في حالة عدد حقيقي عام، وخصوصًا عدد لا نسبي؟}
\lecturerresponse{قد يتطلب العدد الحقيقي الدقيق عددًا لا نهائيًا من البتات. بعض التوسعات الثنائية للأعداد النسبية دورية، والتوسع الثنائي لعدد لا نسبي عام لا ينتهي أبدًا. ويحفظ الحاسوب الحقيقي دائمًا عددًا منتهيًا من الخانات. والمقصود هو أن أي دقة منتهية محددة يمكن ترميزها، لا أن كل عدد حقيقي كلاسيكي له وصف دقيق منتهٍ.}
\end{classroomqa}

\section{حالة وقانون للحركة}

حتى الآن وصفنا نظامًا عند لحظة واحدة. وتتطلب الحركة مكوِّنًا ثانيًا.

\begin{definition}
ال\emph{تهيئة} هي الحالة الكاملة لنظام عند لحظة محددة. وقانون الحركة في الزمن المتقطع هو قاعدة تحوّل كل تهيئة إلى التهيئة بعد خطوة زمنية واحدة.
\end{definition}

هذا التمييز أولي، لكنه جوهري. فاللقطة الآنية لتوزيع إلكترونات هي تهيئة؛ لكنها ليست بعدُ مدارًا. وإذا استُبدل الفضاء بشبكة والزمن بمتتالية \(t_0,t_1,t_2,\ldots\)، فإن الحركة هي التعاقب
\begin{equation}
  \ket{s(t_0)}
  \longmapsto
  \ket{s(t_1)}
  \longmapsto
  \ket{s(t_2)}
  \longmapsto\cdots.
  \label{eq:l01-discrete-motion}
\end{equation}
ولذلك يحتاج النموذج الفيزيائي إلى كل من مجموعة الحالات الممكنة وقاعدة للتحديث.

\begin{classroomqa}
\lecturetimestamp{00:31:12}
\audiencequestion{كيف يعبّر هذا الوصف بالخلايا عن إلكترون يتحرك في مدار؟}
\lecturerresponse{حدّد أولًا التهيئة عند لحظة فحسب: أي الخلايا مشغولة، مع تثبيت النواة إذا كان ذلك جزءًا من النموذج. ثم اجعل الزمن متقطعًا. وتكون الحركة المدارية هي القاعدة التي تحدّث سلسلة إشغال إلى التي تليها. وقد تكون القاعدة معقدة بسبب ترتيب الخلايا، لكنها تظل تحديثًا للمعلومات.}
\end{classroomqa}

ومن المفيد رسم \emph{فضاء حالات} مجرد. فكل نقطة في هذا الفضاء تدل على تهيئة كاملة، لا على موضع في الفضاء العادي. وتحتوي شبكة إشغال \(10\times10\) على مئة بت، ومن ثم على \(2^{100}\) من التهيئات الممكنة.\editorialnote{تذكر المحاضرة الشفوية بإيجاز \(10^{100}\) أثناء عرض المثال؛ لكن عدد البتات يحدد أن القيمة المقصودة هي \(2^{100}\).} وكل واحدة من تلك التهيئات هي نقطة واحدة في فضاء الحالات.

بالنسبة إلى بت واحد لا توجد إلا نقطتان. يترك قانون الهوية كلتيهما ثابتة:
\begin{equation}
  0\mapsto0,
  \qquad
  1\mapsto1.
  \label{eq:l01-identity-update}
\end{equation}
والإمكان الآخر القابل للعكس يبدلهما:
\begin{equation}
  0\mapsto1,
  \qquad
  1\mapsto0.
  \label{eq:l01-flip-update}
\end{equation}
هذا قلّاب. أما المثال الفيزيائي الذي يتطور من تلقاء نفسه فليس مباشرًا بالقدر نفسه كمفتاح إنارة يُشغَّل يدويًا، لأن التدخل يوسّع النظام. ويعطي تعاقب الليل والنهار دورة تقريبية من حالتين؛ كما يتناوب البندول بين الطرفين، مع أن حالته الكاملة تشمل جميع المواضع والسرعات الوسيطة.

\begin{classroomqa}
\lecturetimestamp{00:40:38}
\audiencequestion{هل يكون مفتاح إنارة عادي مثالًا على قانون القلّاب؟}
\lecturerresponse{ليس إذا كان شخص ما يقلبه، لأن المفتاح عندئذ لا يتطور بنفسه؛ بل ينتمي ذلك الشخص إلى النظام الديناميكي الموسَّع. وما نبحث عنه هو نظام مغلق تتناوب حالته وفق قانونه هو. ويعطي تعاقب الليل والنهار الفكرة، بينما يذكّرنا البندول بأن المثال الحقيقي يملك عادة أكثر من حالتين.}
\end{classroomqa}

لنظام مكوَّن من بتّين أربع تهيئات. قد يحمل أحد قوانين التحديث التهيئات الأربع كلها في دورة واحدة. وقد يتألف قانون آخر من دورتين منفصلتين من خطوتين. هذه حلقات في الفضاء المنطقي للتهيئات، وليست بالضرورة حلقات في الفضاء الفيزيائي. وتحدد الأسهم القانون: فمهما تكن التهيئة التي يبدأ منها النظام، فإن اتباع الأسهم يخبرنا بتاريخه اللاحق.

\section{القابلية للعكس والحفاظ على التمييزات}

ليست جميع مخططات الأسهم ذات المنزلة الفيزيائية نفسها. انظر في قاعدة تندمج فيها حالتان مختلفتان في الحالة اللاحقة نفسها:
\begin{equation}
  a\mapsto c,
  \qquad
  b\mapsto c.
  \label{eq:l01-merging-update}
\end{equation}
تكون الحالة التالية محددة سواء بدأنا من \(a\) أو من \(b\)، ولكن بعد الوصول إلى \(c\) لا يمكننا تحديد أيُّ ماضٍ قد حدث. لقد زال التمييز بين \(a\) و\(b\).

\begin{figure}[htbp]
\centering
\includegraphics[width=0.82\linewidth]{lecture_01_figure_02.png}
\caption{تبقى عدة مخططات للتحديث على اللوح أثناء مناقشة القابلية للعكس. وتتجاور دورات مغلقة مع نمط اندماج، ويُكتب إلى جانبها المصطلح الكمي \emph{الوحدوية}.}
\lectureframe{00:53:21}
\end{figure}

\begin{figure}[htbp]
\centering
\begin{tikzpicture}[>=Latex,thick,node distance=1.15cm]
  \node[circle,fill,inner sep=1.8pt,label=above:\(1\)] (r1) {};
  \node[circle,fill,inner sep=1.8pt,right=of r1,label=above:\(2\)] (r2) {};
  \node[circle,fill,inner sep=1.8pt,below=of r2,label=below:\(3\)] (r3) {};
  \node[circle,fill,inner sep=1.8pt,left=of r3,label=below:\(4\)] (r4) {};
  \draw[->] (r1) -- (r2);
  \draw[->] (r2) -- (r3);
  \draw[->] (r3) -- (r4);
  \draw[->] (r4) -- (r1);
  \node[below=0.5cm of r4,anchor=north west] {\small reversible cycle};

  \node[circle,fill,inner sep=1.8pt,right=4.2cm of r1,label=above:\(a\)] (a) {};
  \node[circle,fill,inner sep=1.8pt,below=1.25cm of a,label=below:\(b\)] (b) {};
  \node[circle,fill,inner sep=1.8pt,right=1.6cm of a,yshift=-0.625cm,label=above:\(c\)] (c) {};
  \node[circle,fill,inner sep=1.8pt,right=of c,label=above:\(d\)] (d) {};
  \draw[->] (a) -- (c);
  \draw[->] (b) -- (c);
  \draw[->] (c) -- (d);
  \node[below=0.55cm of b,anchor=north west] {\small information-losing merge};
\end{tikzpicture}
\caption{تعيّن القاعدة القابلة للعكس لكل حالة سابقةً واحدة ولاحقةً واحدة. أما قاعدة الاندماج فلها مستقبل وحيد، ولكن لها ماضٍ ملتبس.}
\end{figure}

لا تسمح الديناميكيات الكلاسيكية الأساسية بهذا الفقدان عندما يُحتفَظ بكل درجة من درجات الحرية. وللخاصية المطلوبة شِقّان:
\begin{equation}
  \boxed{\begin{gathered}
    \text{one unique future and one unique past}\\
    \text{for every state}
  \end{gathered}.}
  \label{eq:l01-reversibility}
\end{equation}
ولذلك، في فضاء حالات منتهٍ، يجب أن يكون التحديث تبديلًا. فلكل حالة سهم خارج واحد بالضبط وسهم داخل واحد بالضبط، ومن ثم يتفكك المخطط إلى دورات منفصلة.

\begin{classroomqa}
\lecturetimestamp{00:46:32}
\audiencequestion{ما المختلف في مخطط الاندماج؟}
\lecturerresponse{إنه لا يتألف بالكامل من حلقات. وعلى نحو أعمق، فإن وجود النظام في الحالة المندمجة لا يخبرنا ما الحالة السابقة التي كانت له. يمكننا بناء المستقبل، لكن لا يمكننا إعادة بناء الماضي. لقد فُقد بتٌّ واحد من التمييز التاريخي.}
\end{classroomqa}

هذا المفهوم هو \emph{قابلية التطور الزمني للعكس}، وليس بالضرورة تناظر الانعكاس الزمني الدقيق. فالدورة التي تسير مع عقارب الساعة لها توجّه أمامي محدد. والقانون الأصلي لا يجعل النظام يسير عكس عقارب الساعة، لكن يوجد قانون معكوس يفعل ذلك. ويطرح تناظر الانعكاس الزمني الدقيق سؤالًا إضافيًا: بعد تطبيق عملية العكس المناسبة، هل نحصل على القوانين الديناميكية نفسها؟ أما القابلية للعكس وحدها فلا تسأل إلا عمّا إذا كان يمكن إعادة تشغيل التاريخ إلى الوراء على نحو فريد.

\begin{classroomqa}
\lecturetimestamp{00:49:58}
\audiencequestion{هل فرادة الماضي هي مجرد تناظر الانعكاس الزمني؟}
\lecturerresponse{ليس تمامًا. فقد تشير أسهم الدورة في اتجاه واحد فقط، ولذلك يكون للقانون توجّه. وتعني القابلية للعكس أنه يمكن إنشاء قانون معكوس محدد جيدًا بعكس كل سهم. وهذا لا يقول، في حد ذاته، إن القانون المعكوس مطابق للقانون الأصلي.}
\end{classroomqa}

ولا ينبغي لنا أيضًا أن نصلح مخططًا غير قابل للعكس بوضع سهمين خارجين من الحالة المندمجة. فهذا سيجعل الحركة إلى الوراء ملتبسةً على نحو مختلف: إذ سيحتاج النظام إلى تعليمة احتمالية، ربما تختار أحد الفرعين في \(30\%\) من المرات والآخر في \(70\%\) منها. وما زلنا نبني ديناميكيات كلاسيكية حتمية، بحيث تحدد الحالة الحاضرة الواحدة مستقبلًا واحدًا وماضيًا واحدًا.

كل نظام حقيقي هو في نهاية المطاف كمومي. وأما لماذا تبدو الأجسام الكبيرة كلاسيكية، ولماذا تصبح بدائلها الكمومية عسيرة الرصد، فذلك سؤال لاحق. وسيحافظ المبدأ الكمومي المناظر على اللب المنطقي نفسه. وعلى سبيل المزاح إلى حد ما، قد يسمّي المرء فرادة الماضي والمستقبل \emph{unique-tarity}. لكن اسمها الحقيقي هو \emph{unitarity}: إذ يمكن للتطور الكمومي أن يجري على نحو فريد إلى الأمام أو إلى الخلف.

\section{الفقدان الظاهري للمعلومات والتخشين}

تبدو الخبرة العيانية وكأنها تناقض القابلية للعكس. ضع جميع جزيئات غاز في الزاوية اليسرى من غرفة محكمة الإغلاق ثم أطلقها. وسرعان ما يملأ الغاز الغرفة. وابدأ بدلًا من ذلك من الزاوية اليمنى فتبدو الصورة العيانية النهائية متماثلة. فهل اندمجت الحالتان الابتدائيتان المتميزتان؟

لا يحدث ذلك إلا بعد أن نكون قد طرحنا التفاصيل جانبًا. فالحالة المجهرية تعيّن كل جزيء. ولو أمكن تتبع كل موضع وكل سرعة بدقة كافية، لأمكن بعكس تلك الحركات التمييز بين التحضيرين. وعمليًا يكون الحساب بالغ الصعوبة إلى حد التعذر، وتنمو الدقة المطلوبة سريعًا، لكن العجز العملي ليس محوًا جوهريًا. والعبارة العيانية \emph{air fills the room} تُعرِّف عددًا هائلًا من التهيئات المجهرية المتميزة بوصفها حالةً واحدة مخشنة.

وهنا يدخل القانون الثاني للديناميكا الحرارية. فنحن نفقد القدرة على التمييز بين الحالات التفصيلية لأننا لا نتتبع جميع المتغيرات المجهرية. وتبدو المعلومات كأنها تختفي لأن وصفنا مخشَّن، لا لأن القانون الحتمي على المستوى الدقيق يدمج نقاطًا متميزة.

\begin{classroomqa}
\lecturetimestamp{00:56:11}
\audiencequestion{هل يتطلب تتبّع جميع الجزيئات في غرفة كلاسيكية عددًا لا نهائيًا من البتات؟}
\lecturerresponse{بالنسبة إلى وصف كلاسيكي دقيق، نعم. فالمواضع والسرعات أعداد حقيقية، وتتبع نظام فوضوي ذي جسيمات كثيرة لأزمنة أطول يتطلب دقةً متزايدة على الدوام. وتغيّر ميكانيكا الكم هذه النتيجة: ففي حيز محكم الإغلاق ذي طاقة محدودة، يمكن أن يكون الوصف القابل للتمييز فيزيائيًا متقطعًا بدلًا من أن يكون قائمةً اعتباطية من أعداد حقيقية كلاسيكية دقيقة.}
\end{classroomqa}

القضية ليست أكاديمية فحسب. وقد كان موضع نزاع لعقود ما إذا كان الثقب الأسود يستطيع تدمير المعلومات؛ وكان Stephen Hawking قد جادل في الأصل بأنه يستطيع ذلك. ويتمحور هذا الجدل حول ما إذا كان الاندماج الظاهري مجرد وصف مخشَّن آخر أم انتهاكًا حقيقيًا للوحدوية الكمومية. ولن نحسمه قبل أن تكون ميكانيكا الكم قد زودتنا باللغة الدقيقة، لكنه يكشف إلى أي حد يعتمد الأمر على مخططات الأسهم البسيطة.

\section{الصفوف والأعمدة والجداءات الداخلية}

قبل أن نستبدل البتات الكلاسيكية بالبتات الكمومية، نحتاج إلى قدر يسير من الجبر الخطي. ولأغراضنا الحالية، يمكن التعامل مع المتجه ببساطة بوصفه قائمة مرتبة منتهية من أعداد حقيقية. ويمكن كتابة متجه ذي أربع مركبات أفقياً،
\begin{equation}
  A^{\mathsf T}=(A_1,A_2,A_3,A_4),
  \label{eq:l01-row-vector}
\end{equation}
أو عموديًا،
\begin{equation}
  A=
  \begin{pmatrix}
    A_1\\ A_2\\ A_3\\ A_4
  \end{pmatrix}.
  \label{eq:l01-column-vector}
\end{equation}
ولا تعني كلمة \emph{four-dimensional} إلا أن هناك أربع مركبات. ولا حاجة إلى تصور سهم في الفضاء المعتاد. فالشكل الأفقي متجه صف، والعمودي متجه عمود.

\begin{figure}[htbp]
\centering
\includegraphics[width=0.92\linewidth]{lecture_01_figure_03.png}
\caption{المركبات نفسها مكتوبة مرةً على هيئة صف ومرةً على هيئة عمود. وفي هذه المرحلة هي أعداد حقيقية، لذا لا يدخل فيها أي اقتران مركب.}
\lectureframe{01:03:05}
\end{figure}

\begin{classroomqa}
\lecturetimestamp{01:03:14}
\audiencequestion{أليس المتجه الصفّي هو المرافق الهيرميتي للمتجه العمودي؟}
\lecturerresponse{سيكون هذا هو الوصف الصحيح حين تكون المركبات عقدية. وعندئذٍ يجب أن يقترن النقل بأخذ المرافق العقدي. أما في الوقت الراهن فالمداخل حقيقية، لذا لا نحتاج إلا إلى التمييز بين الترتيب الصفّي والترتيب العمودي.}
\end{classroomqa}

ليكن
\begin{equation}
  B=(B_1,B_2,B_3,B_4)
\end{equation}
صفًّا له عدد المداخل نفسه الذي للعمود \(A\). ويُعرَّف حاصل ضربهما بمطابقة المواضع المناظرة ثم الجمع:
\begin{align}
  BA
  &=
  (B_1,B_2,B_3,B_4)
  \begin{pmatrix}
    A_1\\ A_2\\ A_3\\ A_4
  \end{pmatrix}\\
  &=B_1A_1+B_2A_2+B_3A_3+B_4A_4.
  \label{eq:l01-inner-product}
\end{align}
والناتج عدد، لا متجه آخر ولا مصفوفة. وفي فضاء متجهي مجرد تُسمّى هذه العملية جداءً داخليًا.

\begin{classroomqa}
\lecturetimestamp{01:06:44}
\audiencequestion{هل هذا هو الجداء النقطي؟}
\lecturerresponse{نعم. بالنسبة إلى المتجهات العادية في الفضاء ثلاثي الأبعاد، فهو الجداء النقطي المألوف. ويؤكد اسم \emph{الجداء الداخلي} أن العملية نفسها، أي ضرب صف في عمود، تنطبق في الفضاءات المتجهية ذات أي بُعد منتهٍ.}
\end{classroomqa}

\section{المصفوفات كآلات لتحديث الحالات}

المصفوفة صفيف مربّع من الأعداد يؤثر في متجه فينتج متجهًا آخر. تخيّلها آلةً: نُدخل العمود القديم، ونطبّق قاعدة ثابتة، فنحصل على عمود جديد. في أربعة أبعاد،
\begin{equation}
  M=
  \begin{pmatrix}
    M_{11}&M_{12}&M_{13}&M_{14}\\
    M_{21}&M_{22}&M_{23}&M_{24}\\
    M_{31}&M_{32}&M_{33}&M_{34}\\
    M_{41}&M_{42}&M_{43}&M_{44}
  \end{pmatrix}.
  \label{eq:l01-general-matrix}
\end{equation}
يعيّن المؤشر الأول الصف، ويعيّن الثاني العمود. ويمكن النظر إلى \(M\) على أنه مجموعة من الأعمدة أو مجموعة من الصفوف؛ وسيكون المنظوران كليهما مفيدين.

لحساب \(A'=MA\)، خذ الجداء الداخلي لكل صف من \(M\) مع عمود الإدخال:
\begin{equation}
  A'_i=(MA)_i=\sum_{j=1}^{4}M_{ij}A_j.
  \label{eq:l01-matrix-vector}
\end{equation}
فعلى سبيل المثال،
\begin{align}
  A'_1&=M_{11}A_1+M_{12}A_2+M_{13}A_3+M_{14}A_4,\\
  A'_2&=M_{21}A_1+M_{22}A_2+M_{23}A_3+M_{24}A_4.
\end{align}
وتتبع المدخلتان الثالثة والرابعة النمط نفسه.

\begin{figure}[htbp]
\centering
\includegraphics[width=0.94\linewidth]{lecture_01_figure_04.png}
\caption{تؤثر المصفوفة في متجه عمودي صفًّا بعد صف. وقد بُسِطت مركبتا الخرج الأوليان صراحةً على اللوح.}
\lectureframe{01:12:45}
\end{figure}

\subsection{ترميز تهيئة كلاسيكية}

لنفترض أن لنظام ما أربع حالات ممكنة. مثّلها بالأعمدة ذات الواحد المنفرد
\begin{equation}
  e_1=
  \begin{pmatrix}1\\0\\0\\0\end{pmatrix},
  \quad
  e_2=
  \begin{pmatrix}0\\1\\0\\0\end{pmatrix},
  \quad
  e_3=
  \begin{pmatrix}0\\0\\1\\0\end{pmatrix},
  \quad
  e_4=
  \begin{pmatrix}0\\0\\0\\1\end{pmatrix}.
  \label{eq:l01-one-hot}
\end{equation}
وموضع \(1\) الوحيد يبيّن لنا أي تهيئة كلاسيكية موجودة.

انظر إلى التحديث الدوري
\begin{equation}
  1\longrightarrow2\longrightarrow3\longrightarrow4\longrightarrow1.
  \label{eq:l01-four-cycle}
\end{equation}
ويمثل بـ
\begin{equation}
  P=
  \begin{pmatrix}
    0&0&0&1\\
    1&0&0&0\\
    0&1&0&0\\
    0&0&1&0
  \end{pmatrix}.
  \label{eq:l01-permutation-matrix}
\end{equation}
وفي الواقع،
\begin{equation}
  \begin{aligned}
    Pe_1&=e_2, & Pe_2&=e_3,\\
    Pe_3&=e_4, & Pe_4&=e_1.
  \end{aligned}
\end{equation}
يصحح البناء المرسوم على اللوح أثناء كتابته، أولا بخمس حالات ثم بأربع. وتسجل المعادلة~\eqref{eq:l01-permutation-matrix} القاعدة الدورية المستقرة لا الترتيب الوسيط المتروك.

تؤثر المصفوفة نفسها في عمود اعتباطي:
\begin{equation}
  P
  \begin{pmatrix}A\\B\\C\\D\end{pmatrix}
  =
  \begin{pmatrix}D\\A\\B\\C\end{pmatrix}.
  \label{eq:l01-cyclic-shift}
\end{equation}
في الحالة الكلاسيكية، لا تكون سوى مركبة واحدة مساوية لـ \(1\)، لكن السماح بمدخلات عامة يكشف العملية الخطية الكامنة. وتستخدم التحديثات الكلاسيكية الحتمية مصفوفات خاصة من الأصفار والآحاد فيها \(1\) واحد في كل صف وكل عمود. وستبقي ميكانيكا الكم على التطور بالمصفوفات مع توسيع مجموعة المدخلات المسموح بها إلى حد كبير.

\subsection{التركيب هو ضرب المصفوفات}

إذا كانت خطوة زمنية واحدة هي
\begin{equation}
  V'=MV,
\end{equation}
فإن تطبيقا ثانيا يعطي
\begin{equation}
  V''=MV'=M(MV)=M^2V.
  \label{eq:l01-two-updates}
\end{equation}
ولذلك فليس ضرب المصفوفات مجرد تمرين جبري اعتباطي. بل إنه يركب التحديثات الفيزيائية. وتمثل خمس خطوات زمنية بـ \(M^5\).

لمصفوفتين \(2\times2\)
\begin{equation}
  M=
  \begin{pmatrix}
    M_{11}&M_{12}\\
    M_{21}&M_{22}
  \end{pmatrix},
  \qquad
  N=
  \begin{pmatrix}
    N_{11}&N_{12}\\
    N_{21}&N_{22}
  \end{pmatrix},
\end{equation}
فإن حاصل الضرب \(C=MN\) له المدخلات
\begin{equation}
  C_{ij}=\sum_k M_{ik}N_{kj}.
  \label{eq:l01-matrix-product}
\end{equation}
وعلى وجه الخصوص،
\begin{align}
  C_{11}&=M_{11}N_{11}+M_{12}N_{21},\\
  C_{12}&=M_{11}N_{12}+M_{12}N_{22},\\
  C_{21}&=M_{21}N_{11}+M_{22}N_{21},\\
  C_{22}&=M_{21}N_{12}+M_{22}N_{22}.
\end{align}
كل مدخلة هي حاصل ضرب صف من المصفوفة اليسرى في عمود من المصفوفة اليمنى. وبصورة مكافئة، أثر بـ \(M\) في كل عمود من \(N\) ثم ضع الأعمدة الناتجة جنبا إلى جنب.

\begin{figure}[htbp]
\centering
\includegraphics[width=0.86\linewidth]{lecture_01_figure_05.png}
\caption{ضرب المصفوفات منظما على هيئة تكرار لعملية ضرب المصفوفة--العمود. وسم الصف الأول وأعمدة المصفوفة الثانية لبيان من أين تأتي كل مدخلة من حاصل الضرب.}
\lectureframe{01:26:30}
\end{figure}

يمكن لمتجه صف أيضا أن يؤثر من اليسار. فإذا كان \(B=(B_1,\ldots,B_4)\)، فإن \(BM\) متجه صف آخر بحيث
\begin{equation}
  (BM)_j=\sum_i B_iM_{ij}.
  \label{eq:l01-row-times-matrix}
\end{equation}
ومرة أخرى، القاعدة هي صف في عمود. وتبين الأشكال هيئة الناتج:
\begin{equation}
\begin{array}{c@{\qquad}c}
\text{row}\times\text{column}&=\text{number},\\
\text{matrix}\times\text{column}&=\text{column},\\
\text{row}\times\text{matrix}&=\text{row},\\
\text{matrix}\times\text{matrix}&=\text{matrix}.
\end{array}
\label{eq:l01-shape-rules}
\end{equation}

ويجدر التدرب على هذه العمليات بأمثلة عددية صغيرة من النوع \(2\times2\) و\(3\times3\). ومع قليل من حساب التفاضل، يعد ضرب الصف--العمود أهم عادة رياضية جديدة لازمة لميكانيكا الكم التي سنقبل عليها.

\begin{classroomqa}
\lecturetimestamp{01:31:50}
\audiencequestion{ما القيد الذي تفرضه القابلية للعكس على مصفوفة التحديث؟}
\lecturerresponse{يجب أن تملك معكوسا. فالمصفوفة المعكوسة تنفذ التحديث العكسي وتعيد كل حالة إلى سابقتها الوحيدة. وليست كل مصفوفة قابلة للعكس، لذا يعيد الجبر إنتاج التمييز الذي أجري سابقا بمخططات الأسهم. ففي المثال الدوري، \(P^{-1}=P^3\) و\(P^4=\mathbb{1}\).}
\end{classroomqa}

\section{لماذا نمهد بالجبر قبل الكيوبت؟}

قدمت الساعة الأولى عرضا نوعيا للفيزياء الكلاسيكية الحتمية من حيث البتات والأسهم القابلة للعكس. وتحول الرياضيات المتبقية ذلك العرض إلى صيغة يمكن أن تبقى صالحة عند الانتقال إلى ميكانيكا الكم. فالتهيئة تصبح عمودا؛ وقانون الخطوة الواحدة يصبح مصفوفة؛ والتطور الزمني المتكرر يصبح حاصل ضرب مصفوفات؛ والقابلية للعكس تصبح وجود معكوس.

لم نستخدم بعد أي قاعدة كمومية. وهذا التأخير مقصود. بعد ذلك سنستبدل البديلين الكلاسيكيين \(\ket{0}\) و\(\ket{1}\) بكيوبت ونسأل: ما الحالات الجديدة الممكنة؟ وحينئذ ستكشف متجهات الصف، ومتجهات العمود، والجداءات الداخلية، والمصفوفات التي قدمت هنا لماذا تختلف المعلومات الكمومية على نحو حاد عن سلسلة كلاسيكية من البتات.

\begin{classroomqa}
\lecturetimestamp{01:34:00}
\audiencequestion{لماذا ندخل كل هذا الجبر المصفوفي قبل أن نقول ما الغرض منه؟}
\lecturerresponse{جاءت المسألة النوعية أولا: كيف ينبغي تمثيل الفيزياء الحتمية بالحالات وقوانين التحديث؟ وهذا الجبر هو التهيئة للخطوة التالية. وما إن تدخل الكيوبتات حتى لن تعود هذه العمليات تبدو منفصلة؛ بل ستكون اللغة الطبيعية للحالات الكمومية وتطورها.}
\end{classroomqa}